% GENERATED FILE. Edit canonical lessons, then run npm run generate:books.
% canonical-source-hash: fnv1a64:597622fa1317b42a
% canonical-lessons: IT-C468-gelo, IT-C468-gelare, IT-C468-brina, IT-C468-rugiada, IT-C468-congelare, IT-R468-ripasso-461-468

\chapter{Il gelo}
\label{ch:it-il-gelo}
\hlchaptermodality{468}

\begin{chapteropening}
\textbf{By the end of this chapter:} \emph{I can describe a cold spell: the frost, freezing, the white frost, the dew and the freezer.}

Complete the last lesson of chapter 468: I can describe a cold spell: the frost, freezing, the white frost, the dew and the freezer.
\end{chapteropening}

\section[il gelo]{il gelo — the frost, the icy cold}

\label{lesson:IT-C468-gelo}



\begin{quote}
Before the new one: say the Italian for the hail, then the Italian for to hail.
\end{quote}

\subsection*{You'll want to know: il gelo}
\textbf{il gelo} — \textquotedblleft{}the frost, the icy cold\textquotedblright{}.

\begin{quote}
Con questo gelo resto a casa.
\end{quote}

\begin{quote}
\emph{In this icy cold I am staying at home.}
\end{quote}

Also a chill between people: \emph{è sceso il gelo}.

\begin{cousinweb}
\emph{gelo} is inherited from Latin \emph{gelū} 'frost, chill' (EtymDB).
\end{cousinweb}

\begin{grammarlens}[title={What you've built}]
One more word for telling, explaining and describing at B1.
\end{grammarlens}

\subsection*{Your turn}
\begin{itemize}
\raggedright
  \item \emph{Say it:} \emph{il gelo}
  \item \emph{Say it:} \emph{il gelo}, once more
  \item \emph{Say it:} a sentence of your own with \emph{il gelo}
  \item \emph{Recall:} say the Italian for the hail, then the Italian for to hail, then say \emph{il gelo} again
\end{itemize}

\begin{tcolorbox}[breakable,colback=teal!4,colframe=teal!35!black,title={Before you move on}]
What does il gelo mean? (\textquotedblleft{}The frost, the icy cold\textquotedblright{}.) Say it once more.
\end{tcolorbox}
\section[gelare]{gelare — to freeze}

\label{lesson:IT-C468-gelare}



\begin{quote}
Before the new one: say the Italian for to hail, then the Italian for the frost.
\end{quote}

\subsection*{You'll want to know: gelare}
\textbf{gelare} — \textquotedblleft{}to freeze\textquotedblright{}.

\begin{quote}
Stanotte è gelato il lago.
\end{quote}

\begin{quote}
\emph{Last night the lake froze.}
\end{quote}

\emph{Si gela!} — it is freezing!

\begin{cousinweb}
\emph{gelare} is inherited from Latin \emph{gelāre} 'to freeze' (EtymDB), from \emph{gelū} 'frost'.
\end{cousinweb}

\begin{grammarlens}[title={What you've built}]
One more word for telling, explaining and describing at B1.
\end{grammarlens}

\subsection*{Your turn}
\begin{itemize}
\raggedright
  \item \emph{Say it:} \emph{gelare}
  \item \emph{Say it:} \emph{gelare}, once more
  \item \emph{Say it:} a sentence of your own with \emph{gelare}
  \item \emph{Recall:} say the Italian for to hail, then the Italian for the frost, then say \emph{gelare} again
\end{itemize}

\begin{tcolorbox}[breakable,colback=teal!4,colframe=teal!35!black,title={Before you move on}]
What does gelare mean? (\textquotedblleft{}To freeze\textquotedblright{}.) Say it once more.
\end{tcolorbox}
\section[la brina]{la brina — the white frost, the hoarfrost}

\label{lesson:IT-C468-brina}



\begin{quote}
Before the new one: say the Italian for the frost, then the Italian for to freeze.
\end{quote}

\subsection*{You'll want to know: la brina}
\textbf{la brina} — \textquotedblleft{}the white frost, the hoarfrost\textquotedblright{}.

\begin{quote}
La mattina i campi sono bianchi di brina.
\end{quote}

\begin{quote}
\emph{In the morning the fields are white with frost.}
\end{quote}

Much thinner than snow.

\begin{cousinweb}
\emph{brina} is inherited from Latin \emph{pruīna} 'hoarfrost' (EtymDB).
\end{cousinweb}

\begin{grammarlens}[title={What you've built}]
One more word for telling, explaining and describing at B1.
\end{grammarlens}

\subsection*{Your turn}
\begin{itemize}
\raggedright
  \item \emph{Say it:} \emph{la brina}
  \item \emph{Say it:} \emph{la brina}, once more
  \item \emph{Say it:} a sentence of your own with \emph{la brina}
  \item \emph{Recall:} say the Italian for the frost, then the Italian for to freeze, then say \emph{la brina} again
\end{itemize}

\begin{tcolorbox}[breakable,colback=teal!4,colframe=teal!35!black,title={Before you move on}]
What does la brina mean? (\textquotedblleft{}The white frost, the hoarfrost\textquotedblright{}.) Say it once more.
\end{tcolorbox}
\section[la rugiada]{la rugiada — the dew}

\label{lesson:IT-C468-rugiada}



\begin{quote}
Before the new one: say the Italian for to freeze, then the Italian for the white frost.
\end{quote}

\subsection*{You'll want to know: la rugiada}
\textbf{la rugiada} — \textquotedblleft{}the dew\textquotedblright{}.

\begin{quote}
L'erba è bagnata di rugiada.
\end{quote}

\begin{quote}
\emph{The grass is wet with dew.}
\end{quote}

In the morning, after a clear night.

\begin{cousinweb}
\emph{rugiada} goes back to Latin \emph{rōs, rōris} 'dew' (EtymDB); its shape points to a Gallo-Romance go-between like Old French \emph{rosée}, though that path is not settled.
\end{cousinweb}

\begin{grammarlens}[title={What you've built}]
One more word for telling, explaining and describing at B1.
\end{grammarlens}

\subsection*{Your turn}
\begin{itemize}
\raggedright
  \item \emph{Say it:} \emph{la rugiada}
  \item \emph{Say it:} \emph{la rugiada}, once more
  \item \emph{Say it:} a sentence of your own with \emph{la rugiada}
  \item \emph{Recall:} say the Italian for to freeze, then the Italian for the white frost, then say \emph{la rugiada} again
\end{itemize}

\begin{tcolorbox}[breakable,colback=teal!4,colframe=teal!35!black,title={Before you move on}]
What does la rugiada mean? (\textquotedblleft{}The dew\textquotedblright{}.) Say it once more.
\end{tcolorbox}
\section[congelare]{congelare — to deep-freeze; to freeze (food)}

\label{lesson:IT-C468-congelare}



\begin{quote}
Before the new one: say the Italian for the white frost, then the Italian for the dew.
\end{quote}

\subsection*{You'll want to know: congelare}
\textbf{congelare} — \textquotedblleft{}to deep-freeze; to freeze (food)\textquotedblright{}.

\begin{quote}
Congeliamo il pane per la settimana.
\end{quote}

\begin{quote}
\emph{We freeze the bread for the week.}
\end{quote}

\emph{Il congelatore} — the freezer.

\begin{cousinweb}
\emph{congelare} is a learned borrowing from Latin \emph{congelāre} 'to freeze solid' (EtymDB), from \emph{gelū} 'frost'.
\end{cousinweb}

\begin{grammarlens}[title={What you've built}]
One more word for telling, explaining and describing at B1.
\end{grammarlens}

\subsection*{Your turn}
\begin{itemize}
\raggedright
  \item \emph{Say it:} \emph{congelare}
  \item \emph{Say it:} \emph{congelare}, once more
  \item \emph{Say it:} a sentence of your own with \emph{congelare}
  \item \emph{Recall:} say the Italian for the white frost, then the Italian for the dew, then say \emph{congelare} again
\end{itemize}

\begin{tcolorbox}[breakable,colback=teal!4,colframe=teal!35!black,title={Before you move on}]
What does congelare mean? (\textquotedblleft{}To deep-freeze; to freeze (food)\textquotedblright{}.) Say it once more.
\end{tcolorbox}
\section[(dialogue)]{Ripasso dei capitoli 461–468}

\label{lesson:IT-R468-ripasso-461-468}



\begin{quote}
Say the Italian for the oak and the pine.
\end{quote}

\subsection*{Your turn}
Say these aloud:

\begin{itemize}
\raggedright
  \item the oak, the pine, the fir, to prune, the bark
  \item the summit, the slope, the ravine, the cave, to climb
  \item the rock, to get a tan, the seashell, the tide, to dive
  \item the wild boar, the deer, the hare, the falcon, to graze
  \item the lizard, the viper, the toad, the snail, to slither
  \item the sultry heat, sultry, to get sunburnt, the drought, the breeze
  \item the flash of lightning, to thunder, to flash, the hail, to hail
  \item the frost, to freeze, the white frost, the dew, to deep-freeze
\end{itemize}

\begin{tcolorbox}[breakable,colback=teal!4,colframe=teal!35!black,title={Before you move on}]
The oak? (\textbf{la quercia}.) The pine? (\textbf{il pino}.) The fir? (\textbf{l'abete}.) To prune? (\textbf{potare}.) The bark? (\textbf{la corteccia}.) The summit? (\textbf{la vetta}.) The slope? (\textbf{il pendio}.) The ravine? (\textbf{il burrone}.) The cave? (\textbf{la caverna}.) To climb? (\textbf{arrampicarsi}.) The rock? (\textbf{lo scoglio}.) To get a tan? (\textbf{abbronzarsi}.) The seashell? (\textbf{la conchiglia}.) The tide? (\textbf{la marea}.) To dive? (\textbf{tuffarsi}.) The wild boar? (\textbf{il cinghiale}.) The deer? (\textbf{il cervo}.) The hare? (\textbf{la lepre}.) The falcon? (\textbf{il falco}.) To graze? (\textbf{brucare}.) The lizard? (\textbf{la lucertola}.) The viper? (\textbf{la vipera}.) The toad? (\textbf{il rospo}.) The snail? (\textbf{la lumaca}.) To slither? (\textbf{strisciare}.) The sultry heat? (\textbf{l'afa}.) Sultry? (\textbf{afoso}.) To get sunburnt? (\textbf{scottarsi}.) The drought? (\textbf{la siccità}.) The breeze? (\textbf{la brezza}.) The flash of lightning? (\textbf{il lampo}.) To thunder? (\textbf{tuonare}.) To flash? (\textbf{lampeggiare}.) The hail? (\textbf{la grandine}.) To hail? (\textbf{grandinare}.) The frost? (\textbf{il gelo}.) To freeze? (\textbf{gelare}.) The white frost? (\textbf{la brina}.) The dew? (\textbf{la rugiada}.) To deep-freeze? (\textbf{congelare}.)
\end{tcolorbox}
